\documentclass[../main.tex]{subfiles}
\begin{document}

\section{Invariants}
\label{sec:kpcfixed:invariants}

Fix a triple $q$ and a history. All handles below have the qualified
membership of \zcref{def:kpcfixed-model:keyed-membership}; validity uses
exact enforcement under \zcref{assum:kpcfixed-model:realization}.

\begin{theorem}[Unique source provenance]
\label{thm:kpcfixed-invariants:source-provenance}
Every member of $C_q(x)$ is $r_q(c)$ for exactly one SOURCE handle $c$
with $c\preceq x$. The sets are finite and uniquely determined by the
dependency history. MERGE introduces no record.
\end{theorem}
\begin{proof}
Induct in dependency order. A SOURCE handle $c$ receives only $r_q(c)$;
its fresh identity fixes the unique source, and $c\preceq c$.
At a MERGE every member comes from an input whose successor is the new
handle, extending the asserted path. The finite union preserves provenance
and introduces no identity. These two cases uniquely determine each set.
\end{proof}

\begin{theorem}[Unique continuation]
\label{thm:kpcfixed-invariants:unique-successor}
Every handle has at most one successor, and a consumed handle has exactly
one, itself a handle. For every handle $g$, at every prefix containing
$g$ exactly one $y$ with $g\preceq y$ is live. For a SOURCE handle $c$,
$r_q(c)\in C_q(y)$ exactly when $c\preceq y$, so distinct live handles
have disjoint sets.
\end{theorem}
\begin{proof}
Single consumption gives at most one consuming event; continue-first
makes its first output a handle, and \zcref{prop:kpcfixed-model:single-birth} makes it the only one.
Only the SOURCE creating $c$ introduces $r_q(c)$. In dependency order,
$r_q(c)\in C_q(y)$ exactly when $y=c$ or some input containing $r_q(c)$
has successor $y$. Induction gives exactly $c\preceq y$. The successor
path from any handle is unique and finite and ends at an unconsumed,
hence live, handle; every earlier handle on that path has been consumed.
Two distinct live handles therefore cannot contain the same record.
\end{proof}

\begin{theorem}[Growth and disjoint gathering]
\label{thm:kpcfixed-invariants:set-membership-growth}
If $x\preceq y$, then $C_q(x)\subseteq C_q(y)$, and the sets a MERGE
gathers are pairwise disjoint.
\end{theorem}
\begin{proof}
By \zcref{thm:kpcfixed-invariants:unique-successor}, a record at $x$ has its
SOURCE handle preceding $x$, hence also $y$, so it belongs to $C_q(y)$.
Immediately before a MERGE its consumed handles are distinct and live;
the same theorem makes their sets disjoint.
\end{proof}

\begin{theorem}[Separation of triples]
\label{thm:kpcfixed-invariants:key-isolation}
For distinct triples $q\ne p$, handles and records of $p$ contribute neither inputs
nor new records to the sets $C_q$. A handle of one triple never has a
successor of the other. This is conditional on the realization's
distinct-program and authentication assumptions
(\zcref{assum:kpcfixed-model:realization}).
\end{theorem}
\begin{proof}
Disjoint membership by \zcref{prop:kpcfixed-model:disjoint-membership}
excludes $K_p$ from $I_q(e)$ and prevents a
$p$-handle from continuing a $q$-handle. A SOURCE handle belongs to
only one triple, fixing the triple of its record. Induction on the recurrence
preserves separation. An event consuming handles of both triples would
require its first output in both disjoint classes, contrary to validity.
Separate events can coexist in one history. A difference in any coordinate
is covered by the same distinct-program premise, not an unconditional
claim that a finite program encoding is injective over arbitrary data.
\end{proof}

\end{document}
